\textbf{Seats}
\begin{enumerate}
    \item 5 mins: Intro slides \par
    \item 10 mins: Finish Example 3.4.5 \par
    \textbf{Motivation}
    \item 20 mins: Finish Example 3.4.6 \par
    \textbf{Check piazza}
    \item 35 mins: Finish Example 3.5.1 \par
    \textbf{Check piazza}
    \item 45 mins: Finish Example 3.5.2 \par
    \textbf{Check piazza}
    \item 50 mins: Wrap up \par
\end{enumerate}
\textbf{Seats}

\vspace{10mm}
\textbf{Example 3.4.5}

Picking up from 3.3.3 we have
\begin{align*}
    \int_0^1 x e^x \, {\rm d}x &= [x e^x]_0^1 - \int_0^1 e^x \, {\rm d} x\\
    &= [x e^x - e^x]_0^1 = e - e + 1 = 1
\end{align*}

\vspace{10mm}
\textbf{Example 3.4.6}

Partial fractions
\begin{align*}
    \frac{x+17}{(x+3)(2x-1)} &= \frac{A}{x+3} + \frac{B}{2x-1} \\
    \therefore \; x+17 &= A(2x-1) + B(x+3)
\end{align*}
Set $x = 1/2$, $7B/2 = 35/2$, so $B = 5$.
Set $x = -3$, $-7A = 14$, so $A = -2$. Therefore
\begin{align*}
    \frac{x+17}{(x+3)(2x-1)} &= \frac{-2}{x+3} + \frac{5}{2x-1} \\
\end{align*}
Now integrate
\begin{align*}
    \int \frac{-2}{x+3} \, {\rm d}x &= -2\ln (x+3) + c_1\\
    \int \frac{5}{2x-1} \, {\rm d}x &= \frac{5}{2} \ln(2x-1) + c_2
\end{align*}
So 
\begin{align*}
\int \frac{x+17}{(x+3)(2x-1)} \, {\rm d}x = -2\ln (x+3) + \frac{5}{2} \ln(2x-1) + C
\end{align*}

\vspace{10mm}
\textbf{Example 3.5.1}

First integral
\begin{align*}
    c_0 \int_0^l \int_0^l 1 \, {\rm d} x \, {\rm d} y = c_0 \int_0^l [x]_0^l \, {\rm d} y &= c_0 l \int_0^l 1 \, {\rm d} y\\
    &= c_0 l[y]_0^l = c_0 l^2\\
\end{align*}
Second integral
\begin{align*}
    12 c_0 \int_0^l \int_0^l xy(1-x) \, {\rm d} x \, {\rm d} y &= 2 c_0 \int_0^l y \int_0^l (x-x^2) \, {\rm d} x \, {\rm d} y\\
    &= 12 c_0 \int_0^l y \left[\frac{x^2}{2} - \frac{x^3}{3} \right]_0^l \, {\rm d} y\\
    &= 12 c_0 \left(\frac{l^2}{2} - \frac{l^3}{3}\right) \int_0^l y \, {\rm d} y\\
    &= 12 c_0 \left(\frac{l^2}{2} - \frac{l^3}{3}\right) \left[\frac{y^2}{2}\right]_0^l \\
    &= 12c_0 \frac{l^2}{2} \left(\frac{l^2}{2} - \frac{l^3}{3}\right)
\end{align*}
So $D = Fh \int_0^l \int_0^l c(x,y) \, {\rm d} x \, {\rm d} y $ for $l=1$ [cm], $h=0.04$ [cm], $F=2.5$ [mol/s] and $c_0 = 10^{4}$ mol/cm$^3$. Therefore, integral 1 gives
\begin{align*}
D &= Fhc_0l^2 = 10^{-5} \\
D &= 12Fhc_0 frac{l^2}{2} \left(\frac{l^2}{2} - \frac{l^3}{3}\right) = 10^{-5} 
\end{align*}
So, the same.

\vspace{10mm}
\textbf{Example 3.5.2}

\begin{align*}
    \int_{-\pi}^\pi \int_0^1 (x^2 + 2x - 1) \cos \left(\frac{y}{2} \right)\, {\rm d} x \, {\rm d} y &= \int_{-\pi}^\pi \cos \left(\frac{y}{2} \right) \left[ \frac{x^3}{3} + x^2 - x\right]_0^1  \, {\rm d} y\\
    &= \frac{1}{3}  \int_{-\pi}^\pi \cos \left(\frac{y}{2} \right) \, {\rm d} y\\
    &= \frac{1}{3} \left[2 \sin \left(\frac{y}{2}\right)\right]_{-\pi}^\pi\\
    &= \frac{4}{3}
\end{align*}
Reversing the order gives
\begin{align*}
    \int_0^1 \int_{-\pi}^\pi (x^2 + 2x - 1) \cos \left(\frac{y}{2} \right) \, {\rm d} y \, {\rm d} x  &= \int_0^1 \left[2 \sin \left(\frac{y}{2}\right)\right]_{-\pi}^\pi (x^2 + 2x - 1)  \, {\rm d} x\\
    &= 4 \int_0^1(x^2 + 2x - 1)  \, {\rm d} x\\
    &= 4  \left[ \frac{x^3}{3} + x^2 - x\right]_0^1\\
    &= \frac{4}{3}
\end{align*}